\section{邺与长安之间，建康城下：东西魏与侯景之乱}
\label{sec:eastern-western-wei-houjing}

\eventanchor{NSL-440-589-167}{五三四年孝武帝逃往关中}，宇文泰迎之入长安；高欢则在邺安置孝静帝和百官。次年西魏、东魏分别以长安和邺为中心成立。两位皇帝与两座都城并不意味着北魏旧有资源被平均切开：关中要重新组织军户、流民、屯田和仓储，河北则依靠邺、晋阳、粮仓与地方官员维系前线。国家分裂首先是供给与任命网络被放进不同地理空间的结果。

\subsection{战争不能脱离俘虏、流民与仓储}

五三七年沙苑之战中宇文泰击败高欢，战后西魏安置俘虏和流民；五三八年河桥、洛阳一线战局反复，五三九年两方一度议和，仍未消除关中、河南的竞争。战役年月可靠，兵数、伤亡和行军路线需保留夸饰风险。较稳的结论是，西魏必须以屯田、军户和仓储支撑较小的关中资源，高欢则通过邺、晋阳和河北官员、军队供给维持东魏。所谓胜负并不能替代对每一处道路、粮源与人口如何被调动的追问。

南方和边缘同样不是旁观者。五四一年交州豪族李贲起兵，梁军的讨伐受气候、道路与补给牵制，五四四年李贲称帝。传世汉文史料和后出的越南记忆对这一过程的称呼、视角并不完全一致，不能把其中任何一方当成无疑的唯一记录。五四五年《玉台新咏》的编纂又保存宫体诗和女性题材文本；它照见的是特定文人、宫廷与文本流通，而不是所有女性或地方社会的声音。

\subsection{侯景把北方军队带到建康}

五四七年侯景叛高澄南降梁。梁朝围绕军队的驻地、粮饷和淮南位置同他冲突，表明接纳一支军队从不是单纯的外交姿态。次年侯景渡江反梁、围攻建康；\eventanchor{NSL-440-589-195}{五四九年城陷}，梁武帝被控制而死，寺院、市场、漕运和户籍秩序都受到破坏。围城中的粮食、疫病和军民经验不能由几段宫廷叙事完整复原，但它们确实让建康原有的资源网络失去正常运作。

五五二年王僧辩、陈霸先等收复建康，次年侯景被部下杀害。叛乱核心瓦解并不等于梁朝自动恢复：京口、江州与建康的户籍、仓储、寺院和水运都需要重新接续，江汉、江东各方又争夺谁能代表梁室。五五四年西魏攻取江陵，梁元帝被杀，大批官员、工匠和图书迁徙或散失。江陵的陷落说明南方政权的重建不能只依赖建康一城，也与关中军政和江汉交通直接相连。

\subsection{新朝的门槛}

五五五年陈霸先杀王僧辩、另立萧方智；北齐的南下支持因江南抵抗与补给困难而撤回。五五六年宇文泰死，宇文护接掌西魏军政；五五七年北周建立，陈霸先受禅建陈。三次政权转换的仪式可以按年序辨认，谁在每一步拥有多少选择并不总能由本纪回答。可以确知的是，侯景之乱与东西魏竞争使军人、流民、文书、城市粮道和宗室名义被不断拆分重组；北周和陈不是从稳定秩序中诞生，而是在这种破碎的网络上重新建立朝廷。

\subsection{来源的层次}

东西魏的年序以《魏书》末期帝纪、《周书》与《北史》为主，梁末主要依靠《梁书》《南史》及《资治通鉴》；它们往往从胜者或后继王朝的位置解释政变和战争。江陵图书、官员与工匠的流散，及交州地方社会的反应，尤其提醒我们不能只用都城材料说明整个区域。本节因而仅将沙苑、侯景、建康与江陵写为资源、道路和政权合法性相互牵动的节点，不虚构城内细节或将史书兵额当作可靠统计。


% 年序补线：本节将既有账本中尚未初次落入正文的条目接入相应历史场景。
\subsection{北魏分裂：关中、河北与建康的危机}
\noindent 534年，东魏与西魏关于\eventanchor{NS-LEDGER-017-1}{北魏分裂为东魏西魏的前期动员与资源准备}的记载仍须回到“北魏分裂为东魏西魏的前期动员与资源准备”逐项比读。 《宋书》《南齐书》《梁书》《陈书》帝纪及列传；《魏书》《北齐书》《周书》；《资治通鉴》宋齐梁陈纪保存可比的线索；废立、即位或政权转换主干可核；礼仪、动机和清洗范围常受胜者叙事影响。\par\smallskip
\noindent \eventanchor{NS-LEDGER-017-2}{北魏分裂为东魏西魏的命令、联盟或地方响应}在534年的东魏与西魏材料中留下的是一项异文问题，其所指的核心事项为“北魏分裂为东魏西魏的命令、联盟或地方响应”。 《宋书》《南齐书》《梁书》《陈书》帝纪及列传；《魏书》《北齐书》《周书》；《资治通鉴》宋齐梁陈纪保存可比的线索；废立、即位或政权转换主干可核；礼仪、动机和清洗范围常受胜者叙事影响。\par\smallskip
\noindent 534年的东魏与西魏并未因此多出一场独立事件；\eventanchor{NS-LEDGER-017-3}{北魏分裂为东魏西魏}提示的是“北魏分裂为东魏西魏”的材料分歧。 《宋书》《南齐书》《梁书》《陈书》帝纪及列传；《魏书》《北齐书》《周书》；《资治通鉴》宋齐梁陈纪保存可比的线索；废立、即位或政权转换主干可核；礼仪、动机和清洗范围常受胜者叙事影响。\par\smallskip
\noindent \eventanchor{NS-LEDGER-017-4}{北魏分裂为东魏西魏的执行中的空间与人员处置}使534年的东魏与西魏年序多出一层检验：所要复核的是“北魏分裂为东魏西魏的执行中的空间与人员处置”，而不是以一个定本覆盖全部传述。 《宋书》《南齐书》《梁书》《陈书》帝纪及列传；《魏书》《北齐书》《周书》；《资治通鉴》宋齐梁陈纪保存可比的线索；废立、即位或政权转换主干可核；礼仪、动机和清洗范围常受胜者叙事影响。\par\smallskip
\noindent 534年，东魏与西魏关于\eventanchor{NS-LEDGER-017-5}{北魏分裂为东魏西魏的后续制度或军政调整}的记载仍须回到“北魏分裂为东魏西魏的后续制度或军政调整”逐项比读。 《宋书》《南齐书》《梁书》《陈书》帝纪及列传；《魏书》《北齐书》《周书》；《资治通鉴》宋齐梁陈纪保存可比的线索；废立、即位或政权转换主干可核；礼仪、动机和清洗范围常受胜者叙事影响。\par\smallskip
\noindent \eventanchor{NS-LEDGER-017-6}{北魏分裂为东魏西魏的异文复核}在534年的东魏与西魏材料中留下的是一项异文问题，其所指的核心事项为“北魏分裂为东魏西魏”。 《宋书》《南齐书》《梁书》《陈书》帝纪及列传；《魏书》《北齐书》《周书》；《资治通鉴》宋齐梁陈纪保存可比的线索；废立、即位或政权转换主干可核；礼仪、动机和清洗范围常受胜者叙事影响。\par\smallskip
\noindent 534年的西魏记录下\eventanchor{NSL-440-589-168}{宇文泰迎孝武帝入长安}：宇文泰迎孝武帝入长安，掌握关中军政。这里的压力来自西部军人需借皇帝提升对关陇豪强与州郡的号召力，随后可见关中成为西魏政权核心，宇文氏的辅政地位确立。 《周书·文帝纪》；《魏书·文帝纪》；《北史》《资治通鉴》留下了年序和结果的线索；本纪和列传可互校；政变动机与参与范围常受胜者叙事影响。\par\smallskip
\noindent \eventanchor{NSL-440-589-169}{高欢迁孝静帝与百官至邺}发生于534年的东魏。它并非只是一则名号变换——高欢迁孝静帝与百官至邺，建立东部政权中心；高欢控制河北资源且不愿受关中皇帝影响，从而邺城取代洛阳成为东魏都城，河北官僚与军镇资源集中。 上述轮廓先依《魏书·出帝纪》《孝静帝纪》；《北史》相关纪传；《资治通鉴》，但本纪和列传可互校；政变动机与参与范围常受胜者叙事影响。\par\smallskip
\noindent 535年，东魏关于\eventanchor{NS-LEDGER-018-1}{高欢经营邺城军政的前期动员与资源准备}的记载仍须回到“高欢经营邺城军政的前期动员与资源准备”逐项比读。 《宋书》《南齐书》《梁书》《陈书》帝纪及列传；《魏书》《北齐书》《周书》；《资治通鉴》宋齐梁陈纪保存可比的线索；制度或政令形成可追索；执行地域、对象和持续时间不能由条文直接推定。\par\smallskip
\noindent \eventanchor{NS-LEDGER-018-2}{高欢经营邺城军政的命令、联盟或地方响应}在535年的东魏材料中留下的是一项异文问题，其所指的核心事项为“高欢经营邺城军政的命令、联盟或地方响应”。 《宋书》《南齐书》《梁书》《陈书》帝纪及列传；《魏书》《北齐书》《周书》；《资治通鉴》宋齐梁陈纪保存可比的线索；制度或政令形成可追索；执行地域、对象和持续时间不能由条文直接推定。\par\smallskip
\noindent 535年的东魏并未因此多出一场独立事件；\eventanchor{NS-LEDGER-018-3}{高欢经营邺城军政}提示的是“高欢经营邺城军政”的材料分歧。 《宋书》《南齐书》《梁书》《陈书》帝纪及列传；《魏书》《北齐书》《周书》；《资治通鉴》宋齐梁陈纪保存可比的线索；制度或政令形成可追索；执行地域、对象和持续时间不能由条文直接推定。\par\smallskip
\noindent \eventanchor{NS-LEDGER-018-4}{高欢经营邺城军政的执行中的空间与人员处置}使535年的东魏年序多出一层检验：所要复核的是“高欢经营邺城军政的执行中的空间与人员处置”，而不是以一个定本覆盖全部传述。 《宋书》《南齐书》《梁书》《陈书》帝纪及列传；《魏书》《北齐书》《周书》；《资治通鉴》宋齐梁陈纪保存可比的线索；制度或政令形成可追索；执行地域、对象和持续时间不能由条文直接推定。\par\smallskip
\noindent 535年，东魏关于\eventanchor{NS-LEDGER-018-5}{高欢经营邺城军政的后续制度或军政调整}的记载仍须回到“高欢经营邺城军政的后续制度或军政调整”逐项比读。 《宋书》《南齐书》《梁书》《陈书》帝纪及列传；《魏书》《北齐书》《周书》；《资治通鉴》宋齐梁陈纪保存可比的线索；制度或政令形成可追索；执行地域、对象和持续时间不能由条文直接推定。\par\smallskip
\noindent \eventanchor{NS-LEDGER-018-6}{高欢经营邺城军政的异文复核}在535年的东魏材料中留下的是一项异文问题，其所指的核心事项为“高欢经营邺城军政”。 《宋书》《南齐书》《梁书》《陈书》帝纪及列传；《魏书》《北齐书》《周书》；《资治通鉴》宋齐梁陈纪保存可比的线索；制度或政令形成可追索；执行地域、对象和持续时间不能由条文直接推定。\par\smallskip
\noindent 535年的西魏并未因此多出一场独立事件；\eventanchor{NS-LEDGER-019-1}{宇文泰控制长安的前期动员与资源准备}提示的是“宇文泰控制长安的前期动员与资源准备”的材料分歧。 《宋书》《南齐书》《梁书》《陈书》帝纪及列传；《魏书》《北齐书》《周书》；《资治通鉴》宋齐梁陈纪保存可比的线索；制度或政令形成可追索；执行地域、对象和持续时间不能由条文直接推定。\par\smallskip
\noindent \eventanchor{NS-LEDGER-019-2}{宇文泰控制长安的命令、联盟或地方响应}使535年的西魏年序多出一层检验：所要复核的是“宇文泰控制长安的命令、联盟或地方响应”，而不是以一个定本覆盖全部传述。 《宋书》《南齐书》《梁书》《陈书》帝纪及列传；《魏书》《北齐书》《周书》；《资治通鉴》宋齐梁陈纪保存可比的线索；制度或政令形成可追索；执行地域、对象和持续时间不能由条文直接推定。\par\smallskip
\noindent 535年，西魏关于\eventanchor{NS-LEDGER-019-3}{宇文泰控制长安}的记载仍须回到“宇文泰控制长安”逐项比读。 《宋书》《南齐书》《梁书》《陈书》帝纪及列传；《魏书》《北齐书》《周书》；《资治通鉴》宋齐梁陈纪保存可比的线索；制度或政令形成可追索；执行地域、对象和持续时间不能由条文直接推定。\par\smallskip
\noindent \eventanchor{NS-LEDGER-019-4}{宇文泰控制长安的执行中的空间与人员处置}在535年的西魏材料中留下的是一项异文问题，其所指的核心事项为“宇文泰控制长安的执行中的空间与人员处置”。 《宋书》《南齐书》《梁书》《陈书》帝纪及列传；《魏书》《北齐书》《周书》；《资治通鉴》宋齐梁陈纪保存可比的线索；制度或政令形成可追索；执行地域、对象和持续时间不能由条文直接推定。\par\smallskip
\noindent 535年的西魏并未因此多出一场独立事件；\eventanchor{NS-LEDGER-019-5}{宇文泰控制长安的后续制度或军政调整}提示的是“宇文泰控制长安的后续制度或军政调整”的材料分歧。 《宋书》《南齐书》《梁书》《陈书》帝纪及列传；《魏书》《北齐书》《周书》；《资治通鉴》宋齐梁陈纪保存可比的线索；制度或政令形成可追索；执行地域、对象和持续时间不能由条文直接推定。\par\smallskip
\noindent \eventanchor{NS-LEDGER-019-6}{宇文泰控制长安的异文复核}使535年的西魏年序多出一层检验：所要复核的是“宇文泰控制长安”，而不是以一个定本覆盖全部传述。 《宋书》《南齐书》《梁书》《陈书》帝纪及列传；《魏书》《北齐书》《周书》；《资治通鉴》宋齐梁陈纪保存可比的线索；制度或政令形成可追索；执行地域、对象和持续时间不能由条文直接推定。\par\smallskip
\noindent 535年，西魏的资源、道路或官署在\eventanchor{NSL-440-589-170}{孝武帝死后}处发生位移：孝武帝死后，宇文泰立元宝炬为文帝，西魏正式建国。关中集团需要稳定皇帝名义以组织行政和军事，其后果是西魏以长安为中心重建官僚体系，与东魏长期对峙。 《周书·文帝纪》；《魏书·文帝纪》；《北史》《资治通鉴》可以支撑这一顺序；本纪和列传可互校；政变动机与参与范围常受胜者叙事影响。\par\smallskip
\noindent 在535年，\eventanchor{NSL-440-589-171}{孝静帝在邺即位}把东魏的局面推向“孝静帝在邺即位，东魏正式建国”。这一着力于高欢需要以北魏皇统维持河北州郡与士族的服从，也令东魏政权依赖高欢军府，皇帝权力受到严格限制。 《魏书·出帝纪》《孝静帝纪》；《北史》相关纪传；《资治通鉴》使这一变化可被追踪，不过本纪和列传可互校；政变动机与参与范围常受胜者叙事影响。\par\smallskip
\noindent 536年的西魏记录下\eventanchor{NSL-440-589-172}{宇文泰整顿关中军队与州郡}：宇文泰整顿关中军队与州郡，吸纳流民和地方豪强。这里的压力来自西魏人口与财力弱于东魏，必须重建可动员的人户，随后可见关陇军政网络逐步形成，府兵制的早期条件得到强化。 材料主要是《周书·文帝纪》；《魏书·文帝纪》；《北史》《资治通鉴》，因此必须保留限度：地方活动见地理、墓志或文书线索；精英材料不代表整体社会。\par\smallskip
\noindent \eventanchor{NSL-440-589-173}{高欢在邺调整河北官员与军队供给}发生于536年的东魏。它并非只是一则名号变换——高欢在邺调整河北官员与军队供给，巩固东魏新政权；政权初建需要恢复战乱后的漕运、粮仓和州郡征收，从而东魏以河北农业和城市资源支撑对西魏的军事优势。 对这一段，只能以《魏书·出帝纪》《孝静帝纪》；《北史》相关纪传；《资治通鉴》核其大意：制度条文可核；实际执行地域、户等与负担不可由诏令直接推定。\par\smallskip
\noindent 537年，东魏与西魏的\eventanchor{NSL-440-589-174}{宇文泰在沙苑击败高欢军}可从“宇文泰在沙苑击败高欢军”这一变化进入。高欢试图以兵力优势迅速消灭关中对手，西魏依托地形与机动防御；后来西魏保住关中，东西魏均需转入长期资源战。 《周书·文帝纪》；《魏书·文帝纪》；《北史》《资治通鉴》留下了年序和结果的线索；战事年月可据编年材料核对；兵数、杀伤与路线须保留史书夸饰风险。\par\smallskip
\noindent 537年的西魏，\eventanchor{NSL-440-589-175}{沙苑战后西魏安置俘虏与战时流民}使“沙苑战后西魏安置俘虏与战时流民，修补关中供给”成为实际的选择与代价。胜利并未改变西魏人口和粮食相对不足的处境，于是军队与地方社会的互相依赖加深，府兵与屯戍安排更重要。 上述轮廓先依《周书·文帝纪》；《魏书·文帝纪》；《北史》《资治通鉴》，但地方活动见地理、墓志或文书线索；精英材料不代表整体社会。\par\smallskip
\noindent \eventanchor{NSL-440-589-176}{河桥、洛阳一线再成双方争夺战场}见于538年，不是抽象的王朝标签：河桥、洛阳一线再成双方争夺战场，战局反复。沙苑后双方都试图控制河南交通和象征性的旧都，并使军事行动消耗大量人力，河南州郡反复易手。 《魏书·出帝纪》《孝静帝纪》；《北史》相关纪传；《资治通鉴》可以支撑这一顺序；战事年月可据编年材料核对；兵数、杀伤与路线须保留史书夸饰风险。\par\smallskip
\noindent 沿着538年的年序，西魏在\eventanchor{NSL-440-589-177}{宇文泰经营关中屯田、仓储与军户供给以承受东魏压力}时经历“宇文泰经营关中屯田、仓储与军户供给以承受东魏压力”。此前西魏无法仅凭一次战役抵消东魏的人口财力优势；此后关中军政体制更重视户籍、土地和军队之间的绑定。 《周书·文帝纪》；《魏书·文帝纪》；《北史》《资治通鉴》使这一变化可被追踪，不过制度条文可核；实际执行地域、户等与负担不可由诏令直接推定。\par\smallskip
\noindent 539年，东魏与西魏的资源、道路或官署在\eventanchor{NSL-440-589-178}{东西魏一度议和并交换使节}处发生位移：东西魏一度议和并交换使节，未能消除河南、关中争端。连年战争使双方均面临补给与内部整合压力，其后果是短暂外交为人员流动和情报交换留下空间，边界依旧军事化。 材料主要是《魏书·出帝纪》《孝静帝纪》；《北史》相关纪传；《资治通鉴》，因此必须保留限度：政权互动可据多方传记校核；族称、疆界和战果须避免按后世分类固化。\par\smallskip
\subsection{侯景与东西魏：道路、粮饷与失守的都城}
\noindent 在540年，\eventanchor{NSL-440-589-179}{高欢继续以晋阳为北方军事重镇}把东魏的局面推向“高欢继续以晋阳为北方军事重镇，调动河北资源支援前线”。这一着力于东魏需兼顾西魏、柔然和北方草原压力，也令晋阳军镇地位上升，成为高氏集团的战略后方。 对这一段，只能以《魏书·出帝纪》《孝静帝纪》；《北史》相关纪传；《资治通鉴》核其大意：战事年月可据编年材料核对；兵数、杀伤与路线须保留史书夸饰风险。\par\smallskip
\noindent 541年的梁与交州记录下\eventanchor{NSL-440-589-180}{交州豪族李贲起兵反梁}：交州豪族李贲起兵反梁，地方治理危机显现。这里的压力来自梁朝远南州郡面对地方豪强、赋役与边地行政能力不足，随后可见交州脱离梁朝有效控制，南方海陆边疆出现新政权实验。 《梁书》相关列传；《大越史记全书》后出记载；《资治通鉴》留下了年序和结果的线索；战事年月可据编年材料核对；兵数、杀伤与路线须保留史书夸饰风险。\par\smallskip
\noindent \eventanchor{NSL-440-589-181}{梁军讨伐李贲}发生于542年的梁与交州。它并非只是一则名号变换——梁军讨伐李贲，交州战争与气候、道路和补给困难相互影响；梁试图恢复岭南以南的州郡秩序，但远征距离长，从而中央对交州的控制更显脆弱，地方武装获得持续空间。 上述轮廓先依《梁书》相关列传；《大越史记全书》后出记载；《资治通鉴》，但战事年月可据编年材料核对；兵数、杀伤与路线须保留史书夸饰风险。\par\smallskip
\noindent 543年，东魏与西魏的\eventanchor{NSL-440-589-182}{邙山之战中高欢与宇文泰激战}可从“邙山之战中高欢与宇文泰激战，双方均未取得决定性歼灭”这一变化进入。河南战场牵动两国主力，均希望突破对方政治核心；后来长期均势迫使东西魏持续扩军、征粮与争取地方势力。 《魏书·出帝纪》《孝静帝纪》；《北史》相关纪传；《资治通鉴》可以支撑这一顺序；战事年月可据编年材料核对；兵数、杀伤与路线须保留史书夸饰风险。\par\smallskip
\noindent 544年的万春国，\eventanchor{NSL-440-589-183}{李贲在交州称帝}使“李贲在交州称帝，建立万春国的政治名号”成为实际的选择与代价。梁朝讨伐未能消除其地方支持和军事据点，于是这一政权体现交州本地精英对中原王朝秩序的重新塑造，史料多为后出叙述。 《梁书》相关列传；《大越史记全书》后出记载；《资治通鉴》使这一变化可被追踪，不过本纪和列传可互校；政变动机与参与范围常受胜者叙事影响。\par\smallskip
\noindent \eventanchor{NSL-440-589-184}{宇文泰持续以关中军户和部曲补充军队}见于544年，不是抽象的王朝标签：宇文泰持续以关中军户和部曲补充军队，整合地方豪强。西魏须把有限人户转化为可持续的军事组织，并使关陇军人集团与地方社会的结合成为北周、隋的制度背景。 材料主要是《周书·文帝纪》；《魏书·文帝纪》；《北史》《资治通鉴》，因此必须保留限度：地方活动见地理、墓志或文书线索；精英材料不代表整体社会。\par\smallskip
\noindent 沿着545年的年序，梁在\eventanchor{NSL-440-589-185}{徐陵主持或参与《玉台新咏》的编纂活动}时经历“徐陵主持或参与《玉台新咏》的编纂活动，保存宫体诗与女性题材文本”。此前建康文人圈重视诗歌选本、宫廷文化与家族交游；此后作品编年和作者归属需逐篇核验，但反映南朝文学传播网络。 对这一段，只能以《出三藏记集》《弘明集》及序跋；《梁书·武帝纪》；相关校勘研究核其大意：成书、抄写与所述事态须分开；传本年代和作者归属尚有可讨论处。\par\smallskip
\noindent 545年，东魏的资源、道路或官署在\eventanchor{NSL-440-589-186}{高欢继续通过邺、晋阳两地调度军政}处发生位移：高欢继续通过邺、晋阳两地调度军政，稳定河北内部。邙山后东魏无法速胜，必须避免河北豪强和军将离心，其后果是高氏军府的行政渗透加深，皇帝与朝臣空间更受限制。 《魏书·出帝纪》《孝静帝纪》；《北史》相关纪传；《资治通鉴》留下了年序和结果的线索；本纪和列传可互校；政变动机与参与范围常受胜者叙事影响。\par\smallskip
\noindent 在546年，\eventanchor{NSL-440-589-187}{高欢去世}把东魏的局面推向“高欢去世，高澄继承东魏军政主导权”。这一着力于高氏军府已掌握军队、财政与官员任命渠道，也令东魏权力由高欢家族延续，内部将领和士族须重新调整依附。 上述轮廓先依《魏书·出帝纪》《孝静帝纪》；《北史》相关纪传；《资治通鉴》，但本纪和列传可互校；政变动机与参与范围常受胜者叙事影响。\par\smallskip
\noindent 547年的东魏与梁并未因此多出一场独立事件；\eventanchor{NS-LEDGER-020-1}{侯景叛高欢南投梁的前期动员与资源准备}提示的是“侯景叛高欢南投梁的前期动员与资源准备”的材料分歧。 《宋书》《南齐书》《梁书》《陈书》帝纪及列传；《魏书》《北齐书》《周书》；《资治通鉴》宋齐梁陈纪保存可比的线索；核心战事与大致年序可核；军数、战果、路线和地方承受须以异政权记录及考古材料复核。\par\smallskip
\noindent \eventanchor{NS-LEDGER-020-2}{侯景叛高欢南投梁的命令、联盟或地方响应}使547年的东魏与梁年序多出一层检验：所要复核的是“侯景叛高欢南投梁的命令、联盟或地方响应”，而不是以一个定本覆盖全部传述。 《宋书》《南齐书》《梁书》《陈书》帝纪及列传；《魏书》《北齐书》《周书》；《资治通鉴》宋齐梁陈纪保存可比的线索；核心战事与大致年序可核；军数、战果、路线和地方承受须以异政权记录及考古材料复核。\par\smallskip
\noindent 547年，东魏与梁关于\eventanchor{NS-LEDGER-020-3}{侯景叛高欢南投梁}的记载仍须回到“侯景叛高欢南投梁”逐项比读。 《宋书》《南齐书》《梁书》《陈书》帝纪及列传；《魏书》《北齐书》《周书》；《资治通鉴》宋齐梁陈纪保存可比的线索；核心战事与大致年序可核；军数、战果、路线和地方承受须以异政权记录及考古材料复核。\par\smallskip
\noindent \eventanchor{NS-LEDGER-020-4}{侯景叛高欢南投梁的执行中的空间与人员处置}在547年的东魏与梁材料中留下的是一项异文问题，其所指的核心事项为“侯景叛高欢南投梁的执行中的空间与人员处置”。 《宋书》《南齐书》《梁书》《陈书》帝纪及列传；《魏书》《北齐书》《周书》；《资治通鉴》宋齐梁陈纪保存可比的线索；核心战事与大致年序可核；军数、战果、路线和地方承受须以异政权记录及考古材料复核。\par\smallskip
\noindent 547年的东魏与梁并未因此多出一场独立事件；\eventanchor{NS-LEDGER-020-5}{侯景叛高欢南投梁的后续制度或军政调整}提示的是“侯景叛高欢南投梁的后续制度或军政调整”的材料分歧。 《宋书》《南齐书》《梁书》《陈书》帝纪及列传；《魏书》《北齐书》《周书》；《资治通鉴》宋齐梁陈纪保存可比的线索；核心战事与大致年序可核；军数、战果、路线和地方承受须以异政权记录及考古材料复核。\par\smallskip
\noindent \eventanchor{NS-LEDGER-020-6}{侯景叛高欢南投梁的异文复核}使547年的东魏与梁年序多出一层检验：所要复核的是“侯景叛高欢南投梁”，而不是以一个定本覆盖全部传述。 《宋书》《南齐书》《梁书》《陈书》帝纪及列传；《魏书》《北齐书》《周书》；《资治通鉴》宋齐梁陈纪保存可比的线索；核心战事与大致年序可核；军数、战果、路线和地方承受须以异政权记录及考古材料复核。\par\smallskip
\noindent 547年，东魏与梁的资源、道路或官署在\eventanchor{NSL-440-589-188}{侯景叛高澄后降梁}处发生位移：侯景叛高澄后降梁，带领部众进入淮南。高氏内部权力重组及侯景与高澄冲突激化，其后果是梁朝接纳侯景，获得北方兵力也引入难以控制的军阀。 《魏书·出帝纪》《孝静帝纪》；《北史》相关纪传；《资治通鉴》可以支撑这一顺序；战事年月可据编年材料核对；兵数、杀伤与路线须保留史书夸饰风险。\par\smallskip
\noindent 在547年，\eventanchor{NSL-440-589-189}{杨衒之在东魏时期撰成《洛阳伽蓝记》}把东魏的局面推向“杨衒之在东魏时期撰成《洛阳伽蓝记》，追述北魏洛阳寺院与毁乱”。这一着力于北魏洛阳在战乱和政权迁移后衰败，作者以旧都记忆组织材料，也令该书是城市与佛教史重要证词，但怀旧修辞和成书距离须辨析。 《出三藏记集》《弘明集》及序跋；《梁书·武帝纪》；相关校勘研究使这一变化可被追踪，不过成书、抄写与所述事态须分开；传本年代和作者归属尚有可讨论处。\par\smallskip
\noindent 547年的梁记录下\eventanchor{NSL-440-589-190}{梁朝围绕安置侯景军队、粮饷和驻防地点与其发生冲突}：梁朝围绕安置侯景军队、粮饷和驻防地点与其发生冲突。这里的压力来自侯景部众数量多且来源复杂，江南既有财政难以无条件供养，随后可见梁廷与侯景互信迅速破裂，建康安全受到威胁。 材料主要是《梁书》帝纪及列传；《南史》；《资治通鉴》，因此必须保留限度：制度条文可核；实际执行地域、户等与负担不可由诏令直接推定。\par\smallskip
\noindent \eventanchor{NS-LEDGER-021-1}{侯景渡江反梁的前期动员与资源准备}使548年的梁与侯景年序多出一层检验：所要复核的是“侯景渡江反梁的前期动员与资源准备”，而不是以一个定本覆盖全部传述。 《宋书》《南齐书》《梁书》《陈书》帝纪及列传；《魏书》《北齐书》《周书》；《资治通鉴》宋齐梁陈纪保存可比的线索；核心战事与大致年序可核；军数、战果、路线和地方承受须以异政权记录及考古材料复核。\par\smallskip
\noindent 548年，梁与侯景关于\eventanchor{NS-LEDGER-021-2}{侯景渡江反梁的命令、联盟或地方响应}的记载仍须回到“侯景渡江反梁的命令、联盟或地方响应”逐项比读。 《宋书》《南齐书》《梁书》《陈书》帝纪及列传；《魏书》《北齐书》《周书》；《资治通鉴》宋齐梁陈纪保存可比的线索；核心战事与大致年序可核；军数、战果、路线和地方承受须以异政权记录及考古材料复核。\par\smallskip
\noindent \eventanchor{NS-LEDGER-021-3}{侯景渡江反梁}在548年的梁与侯景材料中留下的是一项异文问题，其所指的核心事项为“侯景渡江反梁”。 《宋书》《南齐书》《梁书》《陈书》帝纪及列传；《魏书》《北齐书》《周书》；《资治通鉴》宋齐梁陈纪保存可比的线索；核心战事与大致年序可核；军数、战果、路线和地方承受须以异政权记录及考古材料复核。\par\smallskip
\noindent 548年的梁与侯景并未因此多出一场独立事件；\eventanchor{NS-LEDGER-021-4}{侯景渡江反梁的执行中的空间与人员处置}提示的是“侯景渡江反梁的执行中的空间与人员处置”的材料分歧。 《宋书》《南齐书》《梁书》《陈书》帝纪及列传；《魏书》《北齐书》《周书》；《资治通鉴》宋齐梁陈纪保存可比的线索；核心战事与大致年序可核；军数、战果、路线和地方承受须以异政权记录及考古材料复核。\par\smallskip
\noindent \eventanchor{NS-LEDGER-021-5}{侯景渡江反梁的后续制度或军政调整}使548年的梁与侯景年序多出一层检验：所要复核的是“侯景渡江反梁的后续制度或军政调整”，而不是以一个定本覆盖全部传述。 《宋书》《南齐书》《梁书》《陈书》帝纪及列传；《魏书》《北齐书》《周书》；《资治通鉴》宋齐梁陈纪保存可比的线索；核心战事与大致年序可核；军数、战果、路线和地方承受须以异政权记录及考古材料复核。\par\smallskip
\noindent 548年，梁与侯景关于\eventanchor{NS-LEDGER-021-6}{侯景渡江反梁的异文复核}的记载仍须回到“侯景渡江反梁”逐项比读。 《宋书》《南齐书》《梁书》《陈书》帝纪及列传；《魏书》《北齐书》《周书》；《资治通鉴》宋齐梁陈纪保存可比的线索；核心战事与大致年序可核；军数、战果、路线和地方承受须以异政权记录及考古材料复核。\par\smallskip
\noindent 在548年，\eventanchor{NSL-440-589-191}{侯景举兵反梁}把梁的局面推向“侯景举兵反梁，渡江进攻建康”。这一着力于梁廷试图削减侯景兵权和供给，侯景转而争夺都城，也令江南军府、州镇和宗室的反应分裂，京师陷入长期围困。 对这一段，只能以《梁书》帝纪及列传；《南史》；《资治通鉴》核其大意：战事年月可据编年材料核对；兵数、杀伤与路线须保留史书夸饰风险。\par\smallskip
\noindent 548年的梁记录下\eventanchor{NSL-440-589-192}{侯景围建康}：侯景围建康，城内粮食、疫病与军民秩序恶化。这里的压力来自援军协调失败且水陆交通受阻，随后可见围城将宫廷危机转化为城市社会灾难，南朝首都资源网络遭重创。 《梁书》帝纪及列传；《南史》；《资治通鉴》留下了年序和结果的线索；地方活动见地理、墓志或文书线索；精英材料不代表整体社会。\par\smallskip
\noindent \eventanchor{NS-LEDGER-022-1}{侯景攻陷建康的前期动员与资源准备}使549年的梁年序多出一层检验：所要复核的是“侯景攻陷建康的前期动员与资源准备”，而不是以一个定本覆盖全部传述。 《宋书》《南齐书》《梁书》《陈书》帝纪及列传；《魏书》《北齐书》《周书》；《资治通鉴》宋齐梁陈纪保存可比的线索；核心战事与大致年序可核；军数、战果、路线和地方承受须以异政权记录及考古材料复核。\par\smallskip
\noindent 549年，梁关于\eventanchor{NS-LEDGER-022-2}{侯景攻陷建康的命令、联盟或地方响应}的记载仍须回到“侯景攻陷建康的命令、联盟或地方响应”逐项比读。 《宋书》《南齐书》《梁书》《陈书》帝纪及列传；《魏书》《北齐书》《周书》；《资治通鉴》宋齐梁陈纪保存可比的线索；核心战事与大致年序可核；军数、战果、路线和地方承受须以异政权记录及考古材料复核。\par\smallskip
\noindent \eventanchor{NS-LEDGER-022-3}{侯景攻陷建康}在549年的梁材料中留下的是一项异文问题，其所指的核心事项为“侯景攻陷建康”。 《宋书》《南齐书》《梁书》《陈书》帝纪及列传；《魏书》《北齐书》《周书》；《资治通鉴》宋齐梁陈纪保存可比的线索；核心战事与大致年序可核；军数、战果、路线和地方承受须以异政权记录及考古材料复核。\par\smallskip
\noindent 549年的梁并未因此多出一场独立事件；\eventanchor{NS-LEDGER-022-4}{侯景攻陷建康的执行中的空间与人员处置}提示的是“侯景攻陷建康的执行中的空间与人员处置”的材料分歧。 《宋书》《南齐书》《梁书》《陈书》帝纪及列传；《魏书》《北齐书》《周书》；《资治通鉴》宋齐梁陈纪保存可比的线索；核心战事与大致年序可核；军数、战果、路线和地方承受须以异政权记录及考古材料复核。\par\smallskip
\noindent \eventanchor{NS-LEDGER-022-5}{侯景攻陷建康的后续制度或军政调整}使549年的梁年序多出一层检验：所要复核的是“侯景攻陷建康的后续制度或军政调整”，而不是以一个定本覆盖全部传述。 《宋书》《南齐书》《梁书》《陈书》帝纪及列传；《魏书》《北齐书》《周书》；《资治通鉴》宋齐梁陈纪保存可比的线索；核心战事与大致年序可核；军数、战果、路线和地方承受须以异政权记录及考古材料复核。\par\smallskip
\noindent 549年，梁关于\eventanchor{NS-LEDGER-022-6}{侯景攻陷建康的异文复核}的记载仍须回到“侯景攻陷建康”逐项比读。 《宋书》《南齐书》《梁书》《陈书》帝纪及列传；《魏书》《北齐书》《周书》；《资治通鉴》宋齐梁陈纪保存可比的线索；核心战事与大致年序可核；军数、战果、路线和地方承受须以异政权记录及考古材料复核。\par\smallskip
\noindent 在549年，\eventanchor{NSL-440-589-193}{建康陷落}把梁的局面推向“建康陷落，梁武帝被侯景控制并去世”。这一着力于围城耗尽守军与物资，外援未能形成有效合力，也令梁中央权威崩解，地方将领和宗室各自建立据点。 上述轮廓先依《梁书》帝纪及列传；《南史》；《资治通鉴》，但本纪和列传可互校；政变动机与参与范围常受胜者叙事影响。\par\smallskip
\noindent 549年的梁记录下\eventanchor{NSL-440-589-194}{侯景以萧正德、后又以萧纲等为傀儡操控朝廷}：侯景以萧正德、后又以萧纲等为傀儡操控朝廷。这里的压力来自侯景需要借萧氏皇统合法化对建康和州郡的命令，随后可见傀儡政权无法整合各地军府，江南进入多中心战争。 《梁书》帝纪及列传；《南史》；《资治通鉴》可以支撑这一顺序；本纪和列传可互校；政变动机与参与范围常受胜者叙事影响。\par\smallskip
\subsection{新朝门槛：北齐、北周与陈的接收难题}
\noindent \eventanchor{NSL-440-589-196}{高洋废东魏孝静帝}发生于550年的北齐。它并非只是一则名号变换——高洋废东魏孝静帝，自立为帝，北齐建立；高氏军府长期掌握东魏军队、河北财政与百官任命，从而东魏皇统终结，北齐以邺和晋阳为双重军政中心。 《北齐书》帝纪及列传；《北史·齐本纪》；《资治通鉴》使这一变化可被追踪，不过本纪和列传可互校；政变动机与参与范围常受胜者叙事影响。\par\smallskip
\noindent 550年，北齐的\eventanchor{NSL-440-589-197}{北齐接收东魏旧官僚、军户与州郡文书}可从“北齐接收东魏旧官僚、军户与州郡文书，重整河北统治”这一变化进入。建国须将高氏军府的战时控制转化为常规行政；后来河北的人户、仓储与城市成为北齐对外战争的资源基础。 材料主要是《北齐书》帝纪及列传；《北史·齐本纪》；《资治通鉴》，因此必须保留限度：制度条文可核；实际执行地域、户等与负担不可由诏令直接推定。\par\smallskip
\noindent 550年的北齐，\eventanchor{NSL-440-589-198}{高洋即位后加强对鲜卑、汉人军政集团的赏罚与编制}使“高洋即位后加强对鲜卑、汉人军政集团的赏罚与编制”成为实际的选择与代价。新朝需防范东魏旧臣和高氏宗族内部的竞争，于是军事贵族与官僚均被纳入皇帝个人权威下，稳定性有限。 对这一段，只能以《北齐书》帝纪及列传；《北史·齐本纪》；《资治通鉴》核其大意：本纪和列传可互校；政变动机与参与范围常受胜者叙事影响。\par\smallskip
\noindent \eventanchor{NSL-440-589-199}{侯景废简文帝并自立为汉帝}见于550年，不是抽象的王朝标签：侯景废简文帝并自立为汉帝，继续控制建康。侯景傀儡统治无法获得各地萧氏与将领承认，并使江南战争从都城危机扩展为各路军府争夺的内战。 《梁书》帝纪及列传；《陈书》相关列传；《资治通鉴》留下了年序和结果的线索；本纪和列传可互校；政变动机与参与范围常受胜者叙事影响。\par\smallskip
\noindent 沿着约550年的年序，北齐在\eventanchor{NSL-440-589-294}{北齐在邺城重整东魏旧都的宫署、市场与军政居住空间}时经历“北齐在邺城重整东魏旧都的宫署、市场与军政居住空间”。此前新朝须使高氏军府、百官与河北军户共享同一都城运作体系；此后邺城的供给和坊市网络成为北齐政权的物质基础，考古与文献可互证。 上述轮廓先依《北齐书》帝纪及列传；《北史·齐本纪》；《资治通鉴》，但地方活动见地理、墓志或文书线索；精英材料不代表整体社会。\par\smallskip
\noindent 551年，梁的资源、道路或官署在\eventanchor{NSL-440-589-200}{侯景军与湘东王萧绎、王僧辩等势力在江汉、江东对抗}处发生位移：侯景军与湘东王萧绎、王僧辩等势力在江汉、江东对抗。各方均以梁室名义动员，却缺少统一指挥与粮饷，其后果是长江沿线城镇反复易手，漕运和地方社会持续破坏。 《梁书》帝纪及列传；《陈书》相关列传；《资治通鉴》可以支撑这一顺序；战事年月可据编年材料核对；兵数、杀伤与路线须保留史书夸饰风险。\par\smallskip
\noindent 在551年，\eventanchor{NSL-440-589-201}{北齐与柔然、突厥等草原势力就边境和婚姻关系展开交涉}把北齐的局面推向“北齐与柔然、突厥等草原势力就边境和婚姻关系展开交涉”。这一着力于北齐北部边防需应对柔然衰落与突厥兴起的权力真空，也令草原外交成为北齐与西魏、北周竞争的组成部分。 《周书·突厥传》；《北史·突厥传》；《资治通鉴》使这一变化可被追踪，不过政权互动可据多方传记校核；族称、疆界和战果须避免按后世分类固化。\par\smallskip
\noindent 552年的梁并未因此多出一场独立事件；\eventanchor{NS-LEDGER-023-1}{王僧辩陈霸先平侯景的前期动员与资源准备}提示的是“王僧辩陈霸先平侯景的前期动员与资源准备”的材料分歧。 《宋书》《南齐书》《梁书》《陈书》帝纪及列传；《魏书》《北齐书》《周书》；《资治通鉴》宋齐梁陈纪保存可比的线索；核心战事与大致年序可核；军数、战果、路线和地方承受须以异政权记录及考古材料复核。\par\smallskip
\noindent \eventanchor{NS-LEDGER-023-2}{王僧辩陈霸先平侯景的命令、联盟或地方响应}使552年的梁年序多出一层检验：所要复核的是“王僧辩陈霸先平侯景的命令、联盟或地方响应”，而不是以一个定本覆盖全部传述。 《宋书》《南齐书》《梁书》《陈书》帝纪及列传；《魏书》《北齐书》《周书》；《资治通鉴》宋齐梁陈纪保存可比的线索；核心战事与大致年序可核；军数、战果、路线和地方承受须以异政权记录及考古材料复核。\par\smallskip
\noindent 552年，梁关于\eventanchor{NS-LEDGER-023-3}{王僧辩陈霸先平侯景}的记载仍须回到“王僧辩陈霸先平侯景”逐项比读。 《宋书》《南齐书》《梁书》《陈书》帝纪及列传；《魏书》《北齐书》《周书》；《资治通鉴》宋齐梁陈纪保存可比的线索；核心战事与大致年序可核；军数、战果、路线和地方承受须以异政权记录及考古材料复核。\par\smallskip
\noindent \eventanchor{NS-LEDGER-023-4}{王僧辩陈霸先平侯景的执行中的空间与人员处置}在552年的梁材料中留下的是一项异文问题，其所指的核心事项为“王僧辩陈霸先平侯景的执行中的空间与人员处置”。 《宋书》《南齐书》《梁书》《陈书》帝纪及列传；《魏书》《北齐书》《周书》；《资治通鉴》宋齐梁陈纪保存可比的线索；核心战事与大致年序可核；军数、战果、路线和地方承受须以异政权记录及考古材料复核。\par\smallskip
\noindent 552年的梁并未因此多出一场独立事件；\eventanchor{NS-LEDGER-023-5}{王僧辩陈霸先平侯景的后续制度或军政调整}提示的是“王僧辩陈霸先平侯景的后续制度或军政调整”的材料分歧。 《宋书》《南齐书》《梁书》《陈书》帝纪及列传；《魏书》《北齐书》《周书》；《资治通鉴》宋齐梁陈纪保存可比的线索；核心战事与大致年序可核；军数、战果、路线和地方承受须以异政权记录及考古材料复核。\par\smallskip
\noindent \eventanchor{NS-LEDGER-023-6}{王僧辩陈霸先平侯景的异文复核}使552年的梁年序多出一层检验：所要复核的是“王僧辩陈霸先平侯景”，而不是以一个定本覆盖全部传述。 《宋书》《南齐书》《梁书》《陈书》帝纪及列传；《魏书》《北齐书》《周书》；《资治通鉴》宋齐梁陈纪保存可比的线索；核心战事与大致年序可核；军数、战果、路线和地方承受须以异政权记录及考古材料复核。\par\smallskip
\noindent 552年，突厥与柔然的资源、道路或官署在\eventanchor{NSL-440-589-202}{突厥首领土门击败柔然可汗}处发生位移：突厥首领土门击败柔然可汗，突厥汗国兴起。柔然内部矛盾与对铁勒、突厥部众的控制弱化，其后果是草原霸权转换重塑北朝的贡赐、婚姻和边防战略。 材料主要是《周书·突厥传》；《北史·突厥传》；《资治通鉴》，因此必须保留限度：战事年月可据编年材料核对；兵数、杀伤与路线须保留史书夸饰风险。\par\smallskip
\noindent 在552年，\eventanchor{NSL-440-589-203}{土门死后}把突厥的局面推向“土门死后，突厥汗位由科罗等继承，汗国继续向东西扩展”。这一着力于新兴政权需整合阿史那部众和附属铁勒集团，也令突厥成为北齐、北周均无法忽视的北方强邻。 对这一段，只能以《周书·突厥传》；《北史·突厥传》；《资治通鉴》核其大意：政权互动可据多方传记校核；族称、疆界和战果须避免按后世分类固化。\par\smallskip
\noindent 552年的梁记录下\eventanchor{NSL-440-589-204}{王僧辩、陈霸先等在江南击败侯景主力}：王僧辩、陈霸先等在江南击败侯景主力，收复建康。这里的压力来自侯景统治失去地方支持，反侯势力逐渐协调水陆军，随后可见侯景政权崩解，但战后谁代表梁室的争夺随即展开。 《梁书》帝纪及列传；《陈书》相关列传；《资治通鉴》留下了年序和结果的线索；战事年月可据编年材料核对；兵数、杀伤与路线须保留史书夸饰风险。\par\smallskip
\noindent \eventanchor{NSL-440-589-205}{侯景之乱后建康、京口和江州的户籍、仓储与寺院需重新恢复}发生于552年的梁。它并非只是一则名号变换——侯景之乱后建康、京口和江州的户籍、仓储与寺院需重新恢复；多年围城与战乱破坏长江下游的行政和市场网络，从而江南财政重建成为陈朝长期任务，不能仅以宫廷复位衡量恢复。 上述轮廓先依《梁书》帝纪及列传；《陈书》相关列传；《资治通鉴》，但地方活动见地理、墓志或文书线索；精英材料不代表整体社会。\par\smallskip
\noindent 552年，北齐与柔然残部的\eventanchor{NSL-440-589-295}{北齐在北部边地处置柔然覆亡后逃入境内的部众}可从“北齐在北部边地处置柔然覆亡后逃入境内的部众”这一变化进入。突厥击破柔然后，北朝边境出现新的降附和流徙人群；后来安置、招募与防范并行，史书族称不能直接视为稳定身份。 《北齐书》帝纪及列传；《北史·齐本纪》；《资治通鉴》可以支撑这一顺序；政权互动可据多方传记校核；族称、疆界和战果须避免按后世分类固化。\par\smallskip
\noindent 553年的梁，\eventanchor{NSL-440-589-206}{侯景逃亡途中被部下杀害}使“侯景逃亡途中被部下杀害，叛乱核心集团瓦解”成为实际的选择与代价。军事失败、粮饷枯竭和部众离散使侯景失去控制力，于是梁室虽获喘息，王僧辩、陈霸先与宗室之间的权力竞争加剧。 《梁书》帝纪及列传；《陈书》相关列传；《资治通鉴》使这一变化可被追踪，不过本纪和列传可互校；政变动机与参与范围常受胜者叙事影响。\par\smallskip
\noindent \eventanchor{NSL-440-589-207}{北齐与南朝围绕淮南、江北流民和边将展开安置与交涉}见于553年，不是抽象的王朝标签：北齐与南朝围绕淮南、江北流民和边将展开安置与交涉。侯景之乱改变了江北守军、侨民和地方豪强的依附选择，并使边境人口流动成为双方获取情报和兵源的渠道。 材料主要是《北齐书》帝纪及列传；《北史·齐本纪》；《资治通鉴》，因此必须保留限度：政权互动可据多方传记校核；族称、疆界和战果须避免按后世分类固化。\par\smallskip
\noindent 沿着553年的年序，突厥与柔然在\eventanchor{NSL-440-589-208}{突厥继续攻击柔然残部}时经历“突厥继续攻击柔然残部，部分柔然人向西魏、北齐边地求援或流散”。此前柔然失去草原核心牧地和联盟网络；此后北朝接纳或防范残部，族称与归属在史料中呈现流动性。 对这一段，只能以《周书·突厥传》；《北史·突厥传》；《资治通鉴》核其大意：政权互动可据多方传记校核；族称、疆界和战果须避免按后世分类固化。\par\smallskip
\noindent 554年，西魏与梁关于\eventanchor{NS-LEDGER-024-1}{西魏攻取江陵的前期动员与资源准备}的记载仍须回到“西魏攻取江陵的前期动员与资源准备”逐项比读。 《宋书》《南齐书》《梁书》《陈书》帝纪及列传；《魏书》《北齐书》《周书》；《资治通鉴》宋齐梁陈纪保存可比的线索；核心战事与大致年序可核；军数、战果、路线和地方承受须以异政权记录及考古材料复核。\par\smallskip
\noindent \eventanchor{NS-LEDGER-024-2}{西魏攻取江陵的命令、联盟或地方响应}在554年的西魏与梁材料中留下的是一项异文问题，其所指的核心事项为“西魏攻取江陵的命令、联盟或地方响应”。 《宋书》《南齐书》《梁书》《陈书》帝纪及列传；《魏书》《北齐书》《周书》；《资治通鉴》宋齐梁陈纪保存可比的线索；核心战事与大致年序可核；军数、战果、路线和地方承受须以异政权记录及考古材料复核。\par\smallskip
\noindent 554年的西魏与梁并未因此多出一场独立事件；\eventanchor{NS-LEDGER-024-3}{西魏攻取江陵}提示的是“西魏攻取江陵”的材料分歧。 《宋书》《南齐书》《梁书》《陈书》帝纪及列传；《魏书》《北齐书》《周书》；《资治通鉴》宋齐梁陈纪保存可比的线索；核心战事与大致年序可核；军数、战果、路线和地方承受须以异政权记录及考古材料复核。\par\smallskip
\noindent \eventanchor{NS-LEDGER-024-4}{西魏攻取江陵的执行中的空间与人员处置}使554年的西魏与梁年序多出一层检验：所要复核的是“西魏攻取江陵的执行中的空间与人员处置”，而不是以一个定本覆盖全部传述。 《宋书》《南齐书》《梁书》《陈书》帝纪及列传；《魏书》《北齐书》《周书》；《资治通鉴》宋齐梁陈纪保存可比的线索；核心战事与大致年序可核；军数、战果、路线和地方承受须以异政权记录及考古材料复核。\par\smallskip
\noindent 554年，西魏与梁关于\eventanchor{NS-LEDGER-024-5}{西魏攻取江陵的后续制度或军政调整}的记载仍须回到“西魏攻取江陵的后续制度或军政调整”逐项比读。 《宋书》《南齐书》《梁书》《陈书》帝纪及列传；《魏书》《北齐书》《周书》；《资治通鉴》宋齐梁陈纪保存可比的线索；核心战事与大致年序可核；军数、战果、路线和地方承受须以异政权记录及考古材料复核。\par\smallskip
\noindent \eventanchor{NS-LEDGER-024-6}{西魏攻取江陵的异文复核}在554年的西魏与梁材料中留下的是一项异文问题，其所指的核心事项为“西魏攻取江陵”。 《宋书》《南齐书》《梁书》《陈书》帝纪及列传；《魏书》《北齐书》《周书》；《资治通鉴》宋齐梁陈纪保存可比的线索；核心战事与大致年序可核；军数、战果、路线和地方承受须以异政权记录及考古材料复核。\par\smallskip
\noindent \eventanchor{NSL-440-589-209}{西魏攻陷江陵}见于554年，不是抽象的王朝标签：西魏攻陷江陵，杀梁元帝萧绎并掳掠人口物资。西魏希望控制荆州交通、削弱梁朝并取得江汉资源，并使江陵文化和行政中心遭毁，梁室残余分裂为多支政权。 《周书》帝纪及列传；《北史·周本纪》；《资治通鉴》留下了年序和结果的线索；战事年月可据编年材料核对；兵数、杀伤与路线须保留史书夸饰风险。\par\smallskip
\noindent 沿着554年的年序，西魏在\eventanchor{NSL-440-589-210}{西魏在江陵战后扶植萧詧}时经历“西魏在江陵战后扶植萧詧，建立后梁的政治基础”。此前宇文泰需以萧氏名义治理荆襄并牵制陈与北齐；此后后梁成为西魏、北周在南方的附属缓冲政权。 上述轮廓先依《周书》帝纪及列传；《北史·周本纪》；《资治通鉴》，但本纪和列传可互校；政变动机与参与范围常受胜者叙事影响。\par\smallskip
\noindent 554年，北齐的资源、道路或官署在\eventanchor{NSL-440-589-211}{北齐利用梁朝江陵失陷的局势介入江淮与南朝宗室争夺}处发生位移：北齐利用梁朝江陵失陷的局势介入江淮与南朝宗室争夺。梁室分裂使北齐可与不同萧氏集团建立临时关系，其后果是南北外交更具多边性，朝贡与军事支持常互为条件。 《北齐书》帝纪及列传；《北史·齐本纪》；《资治通鉴》可以支撑这一顺序；政权互动可据多方传记校核；族称、疆界和战果须避免按后世分类固化。\par\smallskip
\noindent 在554年，\eventanchor{NSL-440-589-212}{高洋去世}把北齐的局面推向“高洋去世，高殷即位，杨愔等辅政”。这一着力于高洋晚年政治失序，继承人年幼且宗室强盛，也令皇帝、辅臣与高氏诸王的矛盾很快公开化。 《北齐书》帝纪及列传；《北史·齐本纪》；《资治通鉴》使这一变化可被追踪，不过本纪和列传可互校；政变动机与参与范围常受胜者叙事影响。\par\smallskip
\noindent 554年的西魏与梁记录下\eventanchor{NSL-440-589-296}{江陵陷落后}：江陵陷落后，大批梁朝官员、工匠和图书被迁徙或散失。这里的压力来自战争以都城和长江交通节点为目标，战后掳掠成为军队补给的一部分，随后可见江陵的文化与行政资源被重新分配，南北文献传承受到影响。 材料主要是《周书》帝纪及列传；《北史·周本纪》；《资治通鉴》，因此必须保留限度：地方活动见地理、墓志或文书线索；精英材料不代表整体社会。\par\smallskip
\noindent \eventanchor{NSL-440-589-213}{高演发动政变废高殷}发生于555年的北齐。它并非只是一则名号变换——高演发动政变废高殷，自立为孝昭帝；辅政集团试图约束宗室，反而激发高氏内部权力斗争，从而北齐皇统连续受政变冲击，军事贵族对朝廷稳定更具决定性。 对这一段，只能以《北齐书》帝纪及列传；《北史·齐本纪》；《资治通鉴》核其大意：本纪和列传可互校；政变动机与参与范围常受胜者叙事影响。\par\smallskip
\noindent 555年，梁的\eventanchor{NSL-440-589-214}{王僧辩迎北齐扶持的萧渊明入建康}可从“王僧辩迎北齐扶持的萧渊明入建康，试图重建梁朝”这一变化进入。王僧辩需借北齐支持抗衡萧绎旧部与江南军人；后来引入北齐影响引发陈霸先不满，梁室复兴更加依附外力。 《梁书》帝纪及列传；《陈书》相关列传；《资治通鉴》留下了年序和结果的线索；本纪和列传可互校；政变动机与参与范围常受胜者叙事影响。\par\smallskip
\noindent 555年的梁，\eventanchor{NSL-440-589-215}{陈霸先杀王僧辩}使“陈霸先杀王僧辩，控制建康并改立萧方智”成为实际的选择与代价。江南军人反对北齐干预和王僧辩的政治安排，于是陈霸先掌握梁末政局，为建立陈朝奠定军事实力。 上述轮廓先依《梁书》帝纪及列传；《陈书》相关列传；《资治通鉴》，但本纪和列传可互校；政变动机与参与范围常受胜者叙事影响。\par\smallskip
\noindent \eventanchor{NSL-440-589-216}{北齐军南下支持萧渊明}见于555年，不是抽象的王朝标签：北齐军南下支持萧渊明，后因江南抵抗与补给困难撤退。北齐试图把梁朝重建转为自身南方影响力，并使江淮远征成本显示北齐难以长期控制长江以南。 《北齐书》帝纪及列传；《北史·齐本纪》；《资治通鉴》可以支撑这一顺序；战事年月可据编年材料核对；兵数、杀伤与路线须保留史书夸饰风险。\par\smallskip
\noindent 沿着556年的年序，西魏在\eventanchor{NSL-440-589-217}{宇文泰去世}时经历“宇文泰去世，宇文护接掌西魏军政”。此前宇文泰长期掌握关中军队与官僚，继承安排集中于宇文氏家族；此后宇文护成为建立北周并主导皇帝废立的关键人物。 《周书》帝纪及列传；《北史·周本纪》；《资治通鉴》使这一变化可被追踪，不过本纪和列传可互校；政变动机与参与范围常受胜者叙事影响。\par\smallskip
\noindent 556年，北齐的资源、道路或官署在\eventanchor{NSL-440-589-218}{高演去世}处发生位移：高演去世，高湛即位为武成帝。孝昭帝在位短暂，皇位再度转入高氏支系，其后果是北齐内廷的宗室政治未获根本稳定，近臣影响上升。 材料主要是《北齐书》帝纪及列传；《北史·齐本纪》；《资治通鉴》，因此必须保留限度：本纪和列传可互校；政变动机与参与范围常受胜者叙事影响。\par\smallskip
\noindent 在556年，\eventanchor{NSL-440-589-219}{陈霸先整合京口、建康和江东军队}把梁的局面推向“陈霸先整合京口、建康和江东军队，压制梁末宗室竞争者”。这一着力于王僧辩被杀后需迅速控制长江下游的军事节点，也令江南军人集团成为新的建国支柱，旧梁官僚被重新吸纳。 对这一段，只能以《梁书》帝纪及列传；《陈书》相关列传；《资治通鉴》核其大意：战事年月可据编年材料核对；兵数、杀伤与路线须保留史书夸饰风险。\par\smallskip
